\CUSTOMappendix{Написание кеша для метаданных}
\label{appendix:camunda-config}

\begin{lstlisting}[
  caption={pom.xml},
  label={lst:pom.xml}
]
<?xml version="1.0" encoding="UTF-8"?>

<project xmlns="http://maven.apache.org/POM/4.0.0" xmlns:xsi="http://www.w3.org/2001/XMLSchema-instance"
  xsi:schemaLocation="http://maven.apache.org/POM/4.0.0 http://maven.apache.org/xsd/maven-4.0.0.xsd">
  <modelVersion>4.0.0</modelVersion>

  <parent>
    <relativePath>../pom.xml</relativePath>
  </parent>

  <packaging>jar</packaging>

  <dependencies>
    <dependency>
      <groupId>com.fasterxml.jackson.core</groupId>
      <artifactId>jackson-databind</artifactId>
      <scope>provided</scope>
    </dependency>
    <dependency>
        <groupId>org.springframework</groupId>
        <artifactId>spring-context</artifactId>
        <version>7.0.0</version>
    </dependency>
    <dependency>
        <groupId>org.slf4j</groupId>
        <artifactId>slf4j-api</artifactId>
        <version>2.0.13</version>
    </dependency>
    <!-- nuxeo -->
    <dependency>
      <groupId>org.nuxeo.client</groupId>
      <artifactId>nuxeo-java-client</artifactId>
    </dependency>
    <dependency>
      <groupId>org.projectlombok</groupId>
      <artifactId>lombok</artifactId>
      <scope>provided</scope>
    </dependency>
    <dependency>
      <groupId>org.junit.jupiter</groupId>
      <artifactId>junit-jupiter</artifactId>
      <version>5.10.1</version>
      <scope>test</scope>
    </dependency>
    <dependency>
      <groupId>org.mockito</groupId>
      <artifactId>mockito-core</artifactId>
      <version>5.11.0</version>
      <scope>test</scope>
    </dependency>
    <dependency>
      <groupId>org.mockito</groupId>
      <artifactId>mockito-junit-jupiter</artifactId>
      <version>5.11.0</version>
      <scope>test</scope>
    </dependency>
  </dependencies>

  <repositories>
    
  </repositories>

  <build>
    <plugins>
      <plugin>
        <artifactId>maven-compiler-plugin</artifactId>
        <version>3.8.1</version>
        <configuration>
            <source>17</source>
            <target>17</target>
        </configuration>
      </plugin>
      <plugin>
        <artifactId>maven-jar-plugin</artifactId>
        <version>3.2.0</version>
      </plugin>
      <plugin>
        <artifactId>maven-surefire-plugin</artifactId>
        <version>3.0.0</version>
      </plugin>
    </plugins>
  </build>
</project>
\end{lstlisting}

\begin{lstlisting}[
  caption={DictionaryItem},
  label={lst:DictionaryItem}
]
public class DictionaryItem {

  private static final String IS_DEFAULT = "_is_default";
  private static final String VALUE = "_value";
  private static final String PARENT = "_parent";
  private static final String LABEL = "_label";
  private static final String ID = "_id";

  private final Map<String, Object> additionalFields;

  public DictionaryItem() {
    this.additionalFields = new HashMap<>();
  }
  
  public DictionaryItem(Map<String, Object> additionalFields) {
    this.additionalFields = new HashMap<>(additionalFields);
  }

  @JsonAnySetter
  public void setAdditionalField(String key, Object value) {
    additionalFields.put(key, value);
  }

  @JsonAnyGetter
  public Map<String, Object> getAdditionalFields() {
    return additionalFields;
  }

  private <T> T getField(String key, Class<T> retType) {
    Object valueByKey = additionalFields.get(key);
    if (retType.isInstance(valueByKey)) {
      return (T) valueByKey;
    }
    return null;
  }

  public Boolean getIsDefault() {
    return getField(IS_DEFAULT, Boolean.class);
  }

  public String getValue() {
    return getField(VALUE, String.class);
  }

  public String getParent() {
    return getField(PARENT, String.class);
  }
  
  public String getLabel() {
    return getField(LABEL, String.class);
  }

  public String getId() {
    return getField(ID, String.class);
  }

}
\end{lstlisting}

\begin{lstlisting}[
  caption={RdmDataProvider},
  label={lst:RdmDataProvider}
]
public interface RdmDataProvider {

  Object getAllVocabularies(UUID rootUuid)
      throws IOException, InterruptedException;
      
  Object getAllVocabularies()
      throws IOException, InterruptedException;
}
\end{lstlisting}

\begin{lstlisting}[
  caption={DataProviderRdmClient},
  label={lst:DataProviderRdmClient}
]
@Slf4j
@Service
public class DataProviderRdmClient {

  private static final Duration REQUEST_TIMEOUT = Duration.ofMinutes(30);
  private final HttpClient httpClient = HttpClient.newBuilder()
      .connectTimeout(REQUEST_TIMEOUT)
      .build();

  private static final ObjectMapper OBJECT_MAPPER = new ObjectMapper();

  @Value("${url}")
  private String url;
  @Value("${admin.username}")
  private String adminUsername;
  @Value("${admin.password}")
  private String adminPassword;

  public Map<String, List<DictionaryItem>> getHierarhy()
      throws IOException, InterruptedException {
    URI uri = URI.create(url + "/api/v1/");
    HttpRequest request = HttpRequest.newBuilder(uri)
        .GET()
        .timeout(REQUEST_TIMEOUT)
        .header("Authorization", buildBasicAuthHeader())
        .header("Accept", "application/json")
        .build();

    HttpResponse<String> response =
        httpClient.send(request, HttpResponse.BodyHandlers.ofString(StandardCharsets.UTF_8));
    if (isSuccess(response.statusCode())) {
      String body = response.body();
      return OBJECT_MAPPER.readValue(body, new TypeReference<>() {});
    }

    log.error("Unexpected response when requesting vocabularies: status={}, uri={}",
        response.statusCode(), uri);
    throw new IllegalStateException("Failed to fetch vocabularies: HTTP " + response.statusCode());
  }

  private String buildBasicAuthHeader() {
    return "Basic " + Base64.getEncoder().encodeToString(
      (adminUsername + ":" + adminPassword).getBytes(StandardCharsets.UTF_8));
  }

  private boolean isSuccess(int statusCode) {
    return statusCode >= 200 && statusCode < 300;
  }
}
\end{lstlisting}

\begin{lstlisting}[
  caption={NuxeoRdmService},
  label={lst:NuxeoRdmService}
]
@Slf4j
@Service
public class NuxeoRdmService {
  @Autowired
  @Qualifier("Client")
  private NuxeoClient client;

  public UUID getUuidByDocRef(String docRef) throws Exception {
    Document answerDocument = client
        .operation("GetDocument")
        .input(docRef).execute();
    if (answerDocument == null) {
      return null;
    }
    String rootUuid = answerDocument.getId();
    log.info("{} uuid:{}", RdmClientConfig.ROOT_DICTIONARIES_DATA_GROUP, rootUuid);
    return UUID.fromString(rootUuid);
  }
}
\end{lstlisting}

\begin{lstlisting}[
  caption={RdmParserServiceStringImpl},
  label={lst:RdmParserServiceStringImpl}
]
@Service
public class RdmParserServiceStringImpl implements RdmParserService {

  private final ObjectMapper objectMapper = new ObjectMapper();

  @Override
  public Map<String, List<DictionaryItem>> parseVocabularies(Object payload) {
    try {
      String strPayload = String.valueOf(payload);
      if (strPayload == null || strPayload.isBlank()) {
        return Map.of();
      }
      JsonNode root = objectMapper.readTree(strPayload);
      Map<String, List<DictionaryItem>> result = new HashMap<>();
      collectDictionaries(root, result);
      return result;
    } catch (JsonProcessingException e) {
      throw new IllegalStateException("Failed to parse payload", e);
    }
  }

  private void collectDictionaries(JsonNode node, Map<String, List<DictionaryItem>> accumulator) {
    if (node == null) {
      return;
    }
    if (node.isObject()) {
      node.fields().forEachRemaining(entry -> {
        String key = entry.getKey();
        JsonNode value = entry.getValue();
        if (key.endsWith(RdmClientConfig.RDM_TYPE_SUFFIX) && value.isArray()) {
          String dictionaryName = key.substring(0, 
              key.length() - RdmClientConfig.RDM_TYPE_SUFFIX.length());
          if (!dictionaryName.isBlank()) {
            for (JsonNode typeNode : value) {
              JsonNode itemsNode = typeNode.get(dictionaryName);
              if (itemsNode != null && itemsNode.isArray()) {
                List<DictionaryItem> items =
                    accumulator.computeIfAbsent(dictionaryName, unused -> new ArrayList<>());
                for (JsonNode itemNode : itemsNode) {
                  items.add(toDictionaryItem(itemNode));
                }
              }
            }
          }
        }
        collectDictionaries(value, accumulator);
      });
    } else if (node.isArray()) {
      for (JsonNode child : node) {
        collectDictionaries(child, accumulator);
      }
    }
  }

  private DictionaryItem toDictionaryItem(JsonNode itemNode) {
    Map<String, Object> additionalFields =
        objectMapper.convertValue(itemNode, new TypeReference<Map<String, Object>>() {});

    return new DictionaryItem(additionalFields);
  }
}
\end{lstlisting}

\begin{lstlisting}[
  caption={RdmLibService},
  label={lst:RdmLibService}
]
@Slf4j
@Service
public class RdmLibService {

  private final RefreshService refreshService;

  private final ConcurrentHashMap<String, CopyOnWriteArrayList<DictionaryItem>> cache = 
      new ConcurrentHashMap<>();
  private final CountDownLatch cacheInitLatch = new CountDownLatch(1);

  public RdmLibService(RefreshService refreshService) {
    this.refreshService = refreshService;
  }

  public List<DictionaryItem> getDictByDocRef(String docRef) {
    awaitCacheInit();
    if (docRef != null && !docRef.isBlank()) {
      CopyOnWriteArrayList<DictionaryItem> list = cache.get(docRef);
      return list != null ? new ArrayList<>(list) : null;
    }
    log.error("docRef blank");
    return null;
  }

  public Map<String, List<DictionaryItem>> getAllDicts() {
    awaitCacheInit();

    return cache.entrySet().stream()
        .collect(HashMap::new,
            (m, e) -> m.put(e.getKey(), new ArrayList<>(e.getValue())),
            HashMap::putAll);
  }

  public void put(String dictionaryName, DictionaryItem dictionaryItem) {
    awaitCacheInit();
    
    String id = dictionaryItem.getId();
    if (id == null) {
      log.warn("DictionaryItem without _id prop");
      return;
    }

    cache.compute(dictionaryName, (key, existingList) -> {
      CopyOnWriteArrayList<DictionaryItem> list = 
          existingList != null 
            ? new CopyOnWriteArrayList<>(existingList) : new CopyOnWriteArrayList<>();
      
      boolean found = false;
      for (int i = 0; i < list.size(); i++) {
        if (id.equals(list.get(i).getId())) {
          list.set(i, dictionaryItem);
          log.info("Update line in {} with id {}", dictionaryName, id);
          found = true;
          break;
        }
      }
      
      if (!found) {
        list.add(dictionaryItem);
        log.info("Create new line in {} with id {}", dictionaryName, id);
      }
      
      return list;
    });
  }

  public boolean isCacheInitialized() {
    return cacheInitLatch.getCount() == 0;
  }

  private void awaitCacheInit() {
    try {
      if (!cacheInitLatch.await(10, TimeUnit.MINUTES)) {
        log.error("Cache initialization timeout");
        throw new IllegalStateException("Cache not initialized within timeout");
      }
    } catch (InterruptedException e) {
      Thread.currentThread().interrupt();
      throw new IllegalStateException("Interrupted while waiting for cache init", e);
    }
  }

  public boolean refreshCache() {
    try {
      log.info("{}:{}", RdmClientConfig.LIB_NAME,
          getClass().getPackage().getImplementationVersion());
          
      Map<String, List<DictionaryItem>> vocabularies = refreshService.refresh();
      
      cache.clear();
      vocabularies.forEach((key, items) -> 
          cache.put(key, new CopyOnWriteArrayList<>(items)));
      
      log.info("Vocabularies upload size: {}", vocabularies.size());
      cacheInitLatch.countDown();
      return true;
    } catch (Exception e) {
      log.error("refreshCache failed", e);
      return false;
    }
  }
}
\end{lstlisting}

\begin{lstlisting}[
  caption={RdmRefreshEvent},
  label={lst:RdmRefreshEvent}
]
@AllArgsConstructor
@Getter
public class RdmRefreshEvent {
  private final String message;
}
\end{lstlisting}

\begin{lstlisting}[
  caption={RefreshEventListener},
  label={lst:RefreshEventListener}
]
@Component
@RequiredArgsConstructor
@Slf4j
@EnableAsync
public class RefreshEventListener {

  private final RdmLibService rdmLibService;

  @Async
  @EventListener
  public void onApplicationEvent(RdmRefreshEvent event) {
    log.info("Refresh event: {}", event.getMessage());
    while (true) {
      Boolean result = rdmLibService.refreshCache();
      if (Boolean.TRUE.equals(result)) {
        log.info("Init success");
        break;
      } else {
        log.warn("Init fail (retry after 1 minute)");
      }
      try {
        Thread.sleep(RdmClientConfig.DELAY_BETWEEN_ATTEMPTS);
      } catch (InterruptedException ie) {
        Thread.currentThread().interrupt();
        log.error("Thread interrupted during sleep: {}", ie.getMessage());
        break;
      }
    }
  }
}
\end{lstlisting}

\begin{lstlisting}[
  caption={NuxeoRdmConfig},
  label={lst:NuxeoRdmConfig}
]
@Configuration
public class NuxeoRdmConfig {

  @Value("${username}")
  private String nuxeoAdminUsername;
  @Value("${password}")
  private String nuxeoAdminPassword;
  @Value("${url}")
  private String nuxeoUrl;
  

  @Bean("Client")
  public NuxeoClient nuxeoRdmClient() {
    return new NuxeoClient.Builder()
        .url(nuxeoUrl)
        .authentication(nuxeoAdminUsername, nuxeoAdminPassword)
        .schemas("*")
        .connect();
  }
}
\end{lstlisting}

\begin{lstlisting}[
  caption={RdmLibServiceTest},
  label={lst:RdmLibServiceTest}
]
@ExtendWith(MockitoExtension.class)
class RdmLibServiceTest {

  @Mock
  private RefreshService refreshService;

  private RdmLibService service;

  @BeforeEach
  void setUp() {
    service = new RdmLibService(refreshService);
  }

  @Test
  void shouldParseMinimalDictionaryStructure() throws IOException {
    Map<String, List<DictionaryItem>> result = parsePayload("single_dictionary.json");

    List<DictionaryItem> items = result.get("MinimalDict");
    assertNotNull(items, "Minimal dictionary should be parsed");
    assertEquals(1, items.size(), "Expected single entry");

    DictionaryItem item = items.get(0);
    assertEquals(Boolean.TRUE, item.getIsDefault());
    assertEquals("value1", item.getValue());
    assertEquals("parent1", item.getParent());
    assertEquals("Bar", item.getAdditionalFields().get("Foo"));
    assertTrue(item.getAdditionalFields().containsKey("_value"));
  }

  @Test
  void shouldParseNullDictionaryStructure() throws IOException {
    Map<String, List<DictionaryItem>> result = parsePayload("null_dictionary.json");

    List<DictionaryItem> items = result.get("NullDict");
    assertNotNull(items, "Minimal dictionary should be parsed");
    assertEquals(1, items.size(), "Expected single entry");

    assertFalse(result.containsKey("_rdm_type"), "Expected to ignore empty postfix key");
    assertFalse(result.containsKey(""), "Expected to ignore empty key");
    assertFalse(result.containsKey("second_rdm_type"), "Expected to ignore unrelated key");

    DictionaryItem item = items.get(0);
    assertEquals(null, item.getIsDefault());
    assertEquals(null, item.getValue());
    assertEquals(null, item.getParent());
    assertEquals("Bar", item.getAdditionalFields().get("Foo"));
    assertTrue(item.getAdditionalFields().containsKey("_value"));
  }

  @Test
  void initShouldFetchVocabularies() throws Exception {
    Map<String, List<DictionaryItem>> parsedData = parsePayload("single_dictionary.json");

    when(refreshService.refresh()).thenReturn(parsedData);

    boolean success = service.refreshCache();

    assertTrue(success, "Service should complete initialization");
    verify(refreshService).refresh();

    List<DictionaryItem> cachedDictionary = service.getDictByDocRef("MinimalDict");
    assertNotNull(cachedDictionary, "Cached dictionary should be present");
    assertFalse(cachedDictionary.isEmpty(), "Cached dictionary should not be empty");
  }

  private Map<String, List<DictionaryItem>> parsePayload(String resourceName)
      throws IOException {
    RdmParserService parserService = new RdmParserServiceStringImpl();
    return parserService.parseVocabularies(readPayload(resourceName));
  }

  private String readPayload(String resourceName) throws IOException {
    ClassLoader classLoader = Thread.currentThread().getContextClassLoader();
    try (InputStream inputStream = classLoader.getResourceAsStream(resourceName)) {
      if (inputStream == null) {
        throw new IOException(resourceName + " resource not found");
      }
      byte[] bytes = inputStream.readAllBytes();
      return new String(bytes, StandardCharsets.UTF_8);
    }
  }
}
\end{lstlisting}

\begin{lstlisting}[
  caption={null dictionary},
  label={lst:null dictionary}
]
{
  "NullDict_rdm_type": [
    {
      "NullDict": [
        {
          "_value": null,
          "_parent": null,
          "Foo": "Bar"
        }
      ],
      "": [],
      "_rdm_type": [
        {
          "_value": null,
          "_parent": null,
          "Foo": "Bar"
        }
      ],
      "second_rdm_type": "ffffffff",
      "tr_rdm_type": [
        null
      ],
      "for_rdm_type": [
        "sadasdasdasd"
      ]
    }
  ]
}
\end{lstlisting}

\begin{lstlisting}[
  caption={single dictionary},
  label={lst:single dictionary}
]
{
  "MinimalDict_rdm_type": [
    {
      "MinimalDict": [
        {
          "_is_default": true,
          "_value": "value1",
          "_parent": "parent1",
          "Foo": "Bar"
        }
      ]
    }
  ]
}
\end{lstlisting}